\chapter{Неоднородная конфигурация наименьшего действия}

Трилемма «корневая голономия --- масса --- конденсат» оставляет
единственную локальную возможность: поле майорановского спаривания должно
быть неоднородным. Проверим, возникает ли такая стационарная конфигурация
автоматически и что именно в ней фиксирует геометрия.

\section{Минимальный функционал Гинзбурга--Ландау}

На generator \(y\) длины \(L\) рассмотрим
\begin{equation}
 F[\Phi]=\int_0^L
 \left[
 |D_y\Phi|^2+
 \frac{\lambda}{2}\left(|\Phi|^2-v^2\right)^2
 \right]dy,
\end{equation}
где \(\oint_y a=\pi/2\), а заряд поля спаривания равен \(-2\).
Для периодического скалярного поля с числом намотки
\(n\in\mathbb Z\) ковариантный импульс равен
\begin{equation}
 k_n=\frac{2\pi n+\pi}{L}.
\end{equation}
Следовательно, геометрия даёт полуцелую решётку импульсов. Две нижние ветви
\[
 n=0,\qquad n=-1
\]
имеют \(k=\pm\pi/L\) и точно вырождены.

\section{Точный порог конденсации}

После минимизации по постоянной амплитуде \(\rho=|\Phi|\) получаем
\begin{equation}
 \rho^2=v^2-\frac{k_n^2}{\lambda}.
\end{equation}
Ненулевая стационарная конфигурация существует только при
\begin{equation}
 \boxed{\lambda v^2>k_n^2}.
\end{equation}
Для нижней ветви
\[
 k_{\min}^2=\frac{\pi^2}{L^2}.
\]
Для единичной \(\mathbb{RP}^3\) систолическая длина равна \(L=\pi\),
поэтому порог
принимает особенно простой вид
\begin{equation}
 \lambda v^2>1.
\end{equation}

Плотность энергии конденсированной ветви и выигрыш относительно
\(\Phi=0\) равны
\begin{align}
 f_{\mathrm{cond}}
 &=k_{\min}^2v^2-\frac{k_{\min}^4}{2\lambda},\\
 f_0-f_{\mathrm{cond}}
 &=\frac{(\lambda v^2-k_{\min}^2)^2}{2\lambda}.
\end{align}
Радиальный гессиан
\begin{equation}
 H_\rho=4(\lambda v^2-k_{\min}^2)
\end{equation}
положителен строго выше порога.

\section{Что фиксировано и что остаётся свободным}

Топология фиксирует полуцелый сдвиг, сопряжённую пару \(n=0,-1\) и
единичную поперечную намотку вокруг сердцевины при наличии конденсата.
Но она не фиксирует \(\lambda v^2\), поэтому не гарантирует сам факт
конденсации. Она также не выбирает знак продольного импульса.

\begin{proposition}[Условная неоднородная конфигурация]
Корневая голономия порождает устойчивую неоднородную конфигурацию
спаривания, если \(\lambda v^2>\pi^2/L^2\). Она возникает в виде двух
вырожденных сопряжённых ветвей. Следовательно, топология определяет
спектральный порог и структуру ветвей, но не масштаб вакуума и не
ориентацию.
\end{proposition}

Поперечный индекс по модулю два условно остаётся равным единице: если
стационарная конфигурация существует, меридиональный поток \(\pi\)
по-прежнему заставляет намотку поля спаривания быть нечётной. Однако
одномерное ядро оператора Боголюбова--де Жена должно быть пересчитано на
действительном фоне наименьшего действия.

\begin{protocol}[Следующая проверка коэффициентов исходного действия]
\begin{enumerate}
 \item вывести \(\lambda v^2\) из спектрального действия до обращения к
 нейтринным данным;
 \item проверить знак \(\lambda v^2-\pi^2/L^2\);
 \item вывести член, расщепляющий ориентации \(n=0,-1\);
 \item решить поперечную задачу Боголюбова--де Жена на выбранном фоне;
 \item вычислить отношение определителей к нормальной ветви.
\end{enumerate}
\end{protocol}

Машинный протокол сохранён в
\texttt{s2t\_nonuniform\_pairing\_saddle\_results.json}.